\subsection{Gaussian transitions at curvature boundaries}
\label{sec:mesoscopic-growing-jets}

The singular metric supported at the center enforces the growing vanishing order in the global correction. Its smooth part is flat inside the quadrature domain, so a bounded function at scale $R$ must restore strict positivity before the $L^2$ estimate is applied.

\Needspace{20\baselineskip}
\begin{lemma}\label{lem:shrinking-order-correction}
For all sufficiently large $m$ there is a uniformly bounded smooth function $b_R$ such that the metric on $L^m\otimes K_C^{-1}$ induced by the singular potential $\psi_m=m\phi_t+b_R$ and the inverse canonical metric of $\nu$ has curvature at least $(c/R^2)\nu$ globally. On exterior annuli of radii
comparable to $R$ the lower bound is $(cq/R^2)\nu$. If an $L^m$-valued
$\bar\partial$ source $\alpha$ is supported on such an exterior annulus,
it has a linear solution $v$ with
\begin{equation}\label{eq:order-preserving-correction}
 \|v\|_{h^m}^2\leq C\frac{R^2}{q}\|\alpha\|_{h^m,\nu}^2.
\end{equation}
If the source vanishes near $x$, $v$ is holomorphic there and vanishes
to order at least $q$.
\end{lemma}
\begin{proof}
Take $b_R=|z|^2/R^2$ on a disk containing the hole, and cut it off before
radius $3R$, with its negative curvature supported outside the hole.
The bump is uniformly bounded and its Hessian is $O(R^{-2})$. On that
support the curvature $m\rho$ is comparable to $q/R^2$, so $q\to\infty$
absorbs the negative part. Inside the hole the bump supplies $c/R^2$.
Elsewhere in a fixed coordinate disk the original curvature is bounded
below by a fixed multiple of $m|z|^{2r-2}$; off that disk it is at least
$cm$. Since $mR^2\to\infty$, all these bounds dominate $c/R^2$ and the
bounded Ricci term of $\nu$.

Apply the pointwise inverse-curvature $L^2$ theorem for singular metrics,
in the formulation recorded in \cref{sec:history-sections}.
For sections use $L^m\otimes K_C^{-1}$ in bidegrees $(1,0)$ and $(1,1)$;
this includes the Ricci term and preserves the original form norms under
the canonical identification. More explicitly, if $H$ is the singular metric on $L^m$ and the twisted curvature is at least $\lambda\nu$ for a continuous $\lambda>0$, the estimate is
\[
 \int_C|v|_H^2\,\mathrm{d}\nu
 \leq\int_C\lambda^{-1}|\alpha|_{H,\nu}^2\,\mathrm{d}\nu.
\]
The inverse curvature on the support is
$O(R^2/q)$. There $\phi_t=\phi$, and the bump is bounded. Choosing the minimal solution in the singular weighted Hilbert space makes the correction linear in its source. Elsewhere
$\mathrm{e}^{-m\phi}\leq C \mathrm{e}^{-\psi_m}$. These comparisons yield \eqref{eq:order-preserving-correction}.
To use a continuous positive curvature lower bound, interpolate from
$c/R^2$ inside the hole to $cq/R^2$ across a thin exterior collar.
For the separated cutoff supports below, choose this collar to miss
the source. More generally, for an exterior source accumulating on
$\Gamma_t$, shrink the collar to the boundary. Dominated convergence
in the inverse-curvature integrals and weak compactness in the fixed
weighted Hilbert space give the same estimate; the boundary has area
zero. No discontinuous lower bound is required by the cited theorem.
Approximation of the positive singular metric gives the same estimate
and a weak solution. The distributional Cauchy--Riemann equation makes it holomorphic where the source vanishes.
Near $x$, weighted integrability is against
$|z|^{-2q}$ times a smooth positive factor. A holomorphic term $z^j$ is
integrable precisely when $j\geq q$; thus the solution has the required
vanishing order.
\end{proof}

\Needspace{20\baselineskip}
\begin{theorem}\label{thm:curvature-boundary-jets}\index{Bergman kernel!curvature boundary transition}
For the domains $D_t$ constructed above, let $d_t(z)$ be signed
Euclidean distance from $\Gamma_t$, positive on the exterior, and let
$p_t(z)$ be the nearest boundary point. For every fixed $M>0$, uniformly
in all directions and $|d_t(z)|\leq MR/\sqrt q$,
\begin{equation}\label{eq:curvature-boundary-transition}
 \frac{\pi a(z)B_{m,q}(z)}{m\rho(z)}
 -\Phi\bigl(2\sqrt{m\rho(p_t(z))}\,d_t(z)\bigr)\longrightarrow0.
\end{equation}
If $d_t(z)=o(R)$ and $\sqrt q\,d_t(z)/R$ tends to $-\infty$ or
$+\infty$, the same normalized kernel tends to $0$ or $1$, respectively.
Thus the corrected center and transition length are $\Gamma_t$ and
$1/(2\sqrt{m\rho})$, for the entire sublinear range.
\end{theorem}
\begin{proof}
Define on the unit disk
\[
 \begin{aligned}
 A_R(w)&=G_R(F_R(w))+r\log(F_R(w)/w),\\
 \Psi_R&=Q_R(F_R(w))-2\operatorname{Re} A_R(w)-2r\log|w|.
 \end{aligned}
\]
The logarithm uses the branch near one. Here $A_R=1/2+O(R)$ and
$F_R=w+O(R)$ in every fixed norm used below. The effective potential is
positive inside, with zero value and gradient on the boundary. For
$w=\mathrm{e}^{v+\mathrm{i}\theta}$ its expansion is
\begin{equation}\label{eq:curvature-ring-expansion}
 \Psi_R=(2r^2+O(R))v^2+O(v^3).
\end{equation}
The exact quadratic coefficient is
$2(Q_R)_{\zeta\bar\zeta}(F_R(\mathrm{e}^{\mathrm{i}\theta}))|F_R'(\mathrm{e}^{\mathrm{i}\theta})|^2$.

The maps need not extend holomorphically across the circle. To obtain
genuine exterior charts, truncate their Taylor series at
$K=\lfloor q^{1/16}\rfloor$, denoting them by $F_K,A_K$.
Smooth Fourier coefficient bounds in a sufficiently high fixed norm
make the truncation error $O(q^{-10})$, including the fixed derivatives
needed here. On $|w|\leq1+c/K$ their derivatives remain bounded and
$F_K-w,A_K-1/2=O(R)$: the factor $(1+c/K)^K\leq \mathrm{e}^c$ controls the sums.
In particular $F_K$ is univalent there for small $R$. With
$\delta=c/K$ and
\[
 \Psi_K=Q_R(F_K)-2\operatorname{Re} A_K-2r\log|w|,
\]
value and first derivative errors on the circle are $O(q^{-10})$.
Uniformly for $|v|\leq\delta$,
\begin{equation}\label{eq:polynomial-ring-expansion}
 \Psi_K=(2r^2+O(R)+O(|v|))v^2+O(q^{-10})(1+|v|).
\end{equation}
On the interior the effective potential has uniform positive quadratic
growth near the circle, a positive gap on compact annuli farther inside,
and logarithmic growth at zero.

In $z=RF_K(w)$ a local holomorphic section of order at least $q$ is
uniquely represented, in the original frame, by
\begin{equation}\label{eq:quadrature-section-gauge}
 s=w^q\exp((q/r)A_K(w))f(w),\qquad
 |s|_{h^m}^2=|f(w)|^2\mathrm{e}^{-(q/r)\Psi_K(w)},
\end{equation}
where $f$ is holomorphic at zero and has only nonnegative Taylor powers.
The volume in this chart is $R^2(1+O(R))\mathrm{d}A_w$. Put
$\eta=\sqrt{rq}\,v$. On each fixed $|\eta|\leq L$,
\[
 (q/r)\Psi_K=2\eta^2+o(1).
\]
For an observation $\eta_0$ in a fixed bounded interval, restriction to
this annulus bounds the kernel from above. Its limiting diagonal
monomial norms, for $f(w)=\sum_{j\geq0}b_jw^j$, are
\[
 n_{j,L}=\frac{2\pi R^2}{\sqrt{rq}}I_L(j/\sqrt{rq}),\qquad
 I_L(s)=\int_{-L}^L \mathrm{e}^{2s\eta-2\eta^2}\,\mathrm{d}\eta.
\]
After normalization by $rq/(\pi R^2)$, the limiting upper kernel is
\begin{equation}\label{eq:positive-mode-ring-kernel}
 \frac12\mathrm{e}^{-2\eta_0^2}\int_0^\infty\frac{\mathrm{e}^{2s\eta_0}}{I_L(s)}\,\mathrm{d}s.
\end{equation}
The sum-to-integral passage uses mesh $1/\sqrt{rq}$. To control its tail,
integrate $I_L$ on a fixed exterior subinterval beyond every observation
$\eta_0$; the resulting exponential bound is uniform in the mesh.
Thus Riemann sums and then $L\to\infty$ are justified. Since
$I_\infty(s)=\sqrt{\pi/2}\,\mathrm{e}^{s^2/2}$, \eqref{eq:positive-mode-ring-kernel} tends to
\begin{equation}\label{eq:gaussian-ring-limit}
 \frac1{\sqrt{2\pi}}\int_0^\infty \mathrm{e}^{-(s-2\eta_0)^2/2}\,\mathrm{d}s
 =\Phi(2\eta_0).
\end{equation}

For the lower bound fix $T$ and use \eqref{eq:quadrature-section-gauge} with
$f$ in the span of $w^j$, $0\leq j\leq T\sqrt{rq}$. Let $\|s\|_{\mathrm{trial}}$ denote the Gaussian trial norm below. Cut off with a radial $\chi$ supported through
$1+\delta$, equal to one through $1+\delta/2$.
The cutoff annulus lies outside the true hole. Equations
\eqref{eq:polynomial-ring-expansion} and \eqref{eq:quadrature-section-gauge} bound the squared source norm,
as a quadratic form on the entire span, by
\[
 \|\bar\partial(\chi s)\|_{h^m,\nu}^2\leq C(R\delta)^{-2}\mathrm{e}^{-cq\delta^2}\|s\|_{\mathrm{trial}}^2.
\]
The linear term $2jv$ is absorbed because $\sqrt q\,\delta\to\infty$.
\Cref{lem:shrinking-order-correction} gives a linear correction $v_s$ satisfying
\begin{equation}\label{eq:shrinking-cutoff-correction}
 \|v_s\|_{h^m}^2\leq C(q\delta^2)^{-1}\mathrm{e}^{-cq\delta^2}\|s\|_{\mathrm{trial}}^2.
\end{equation}
The factors depend on $m$ through $R$ and $q$; the factor $R^2$ in the correction estimate cancels the cutoff derivative. This is needed even when $q$ grows slower than every power of $m$.

The trial Gram form converges in relative operator norm to the diagonal
form
\begin{equation}\label{eq:gaussian-ring-gram}
 n_j=\frac{2\pi R^2}{\sqrt{rq}}\sqrt{\frac\pi2}
          \exp(j^2/(2rq)).
\end{equation}
Indeed the weight ratio is uniform on each fixed $\eta$-band. Outside
it, radial domination and angular Parseval bound arbitrary combinations
by diagonal tails $\mathrm{e}^{-c\eta^2+2T|\eta|}$. First pass to the limit on
the fixed band and then enlarge it. Interior compact gaps and the
logarithmic singularity provide the remaining domination. This controls
the whole quadratic form without a dimension factor.

At points of radius comparable to $R$, a local submean estimate gives
\begin{equation}\label{eq:shrinking-point-evaluation}
 |u(z)|_{h^m}^2\leq C(q/R^2)\|u\|_{h^m}^2
\end{equation}
for a section $u$ holomorphic near the point on a disk of radius
$R/\sqrt q$. Remove the harmonic linear and quadratic Taylor terms
there. The scaled remaining quadratic coefficient is bounded and the
third-order error is $O(q^{-1/2})$, which proves the estimate.
It applies to the corrections, since their sources lie on the distant
cutoff annulus. Together with \eqref{eq:shrinking-cutoff-correction}, it shows that the
corrected global trial kernel differs by $o(q/R^2)$ from the local trial
kernel. The normalized lower limit is \eqref{eq:gaussian-ring-limit} with upper
limit $T$ instead of infinity; let $T\to\infty$.

At a nearest boundary point,
\[
 \begin{aligned}
 v&=\frac{d_t(z)}{R|F_R'(\mathrm{e}^{\mathrm{i}\theta})|}
       +O(d_t(z)^2/R^2)+o(q^{-1/2}),\\
 m\rho(p)R^2|F_R'(\mathrm{e}^{\mathrm{i}\theta})|^2&=rq(1+O(R)).
 \end{aligned}
\]
Also $m\rho(z)/a(z)=(rq/R^2)(1+o(1))$ in the fixed transition band.
These identities convert $2\eta_0$ to the argument in \eqref{eq:curvature-boundary-transition}
and give its exact density normalization.

\emph{Full density and saturation.}
At points $|z|\asymp R$ the full, unselected kernel obeys
\begin{equation}\label{eq:mesoscopic-full-density}
 B_m^{\rm full}(z)=\frac{m\rho(z)}{\pi a(z)}(1+o(1)),
\end{equation}
uniformly in an $o(R)$ neighborhood of $\Gamma_t$. Here is a direct
localization argument, which also supplies the saturation part of
\cref{thm:curvature-boundary-jets}. Remove the harmonic Taylor terms of $m\phi$
at the observation and rescale by $(m\rho(z))^{-1/2}$. On every fixed
rescaled disk the remaining weight tends to $|u|^2$ and the volume to
the constant $a(z)$. Taylor's estimate is uniform because the third
derivatives at radius $R$ are $O(R^{2r-3})$, so the normalized third
order error is $O(q^{-1/2})$. Restriction to such disks bounds the
evaluation norm from above; their Gaussian constant-function norm tends
to $\pi$ as their radius grows. This proves the upper bound in
\eqref{eq:mesoscopic-full-density}.

For the lower bound use the holomorphic exponential cancelling the
harmonic Taylor terms, and cut it off on a disk of radius $\varepsilon R$
centered at the observation, where $\varepsilon\to0$ and
$\sqrt q\,\varepsilon\to\infty$. Its local squared norm is
$(\pi a(z)/(m\rho(z)))(1+o(1))$ relative to its unit weighted value
at the center. Quadratic lower domination bounds the source annulus by
an exponential $\mathrm{e}^{-cq\varepsilon^2}$. Add the same bounded bump of
scale $R$ to the original metric and use the pointwise $\bar\partial$
estimate. The relative squared correction is bounded by
$C(q\varepsilon^2)^{-1}\mathrm{e}^{-cq\varepsilon^2}=o(1)$, after decreasing
$c$ if needed. Its value at the center is $o(1)$ by the local submean
estimate \eqref{eq:shrinking-point-evaluation}. This proves the matching lower bound.

For interior saturation, use the polynomial chart above and put
$u=1-|w|>0$. Restrict the norm in \eqref{eq:quadrature-section-gauge} to the exterior
half-band $0\leq v\leq c/\sqrt q$. Every nonnegative Taylor monomial
of $f$ has norm at least $cR^2/\sqrt q$. Angular Parseval and
Cauchy--Schwarz give
\[
 B_{m,q}(z)\leq
 C\frac{\sqrt q}{R^2}\frac{\mathrm{e}^{-cqu^2}}{1-|w|^2}.
\]
If $u\sqrt q\to\infty$, division by $q/R^2$ makes this tend to zero.
The chart error is $O(q^{-10})$ and does not alter the hypothesis or
the conclusion.

For exterior saturation, use a Gaussian peak at $z$ and cut it off on
a disk of radius
\[
 \varepsilon R,\qquad
 \varepsilon=\min\{d_t(z)/(3R),q^{-1/4}\}.
\]
It is wholly outside the hole. Again $\varepsilon\to0$ and
$\sqrt q\varepsilon\to\infty$. Apply \cref{lem:shrinking-order-correction} using the
singular obstacle metric, which equals the original metric on this
source. The correction has negligible norm and value and vanishes to
order $q$ at $x$. Thus a member of $V_{m,q}$ attains the full Gaussian
density asymptotically. Equation~\eqref{eq:mesoscopic-full-density} is its upper bound.
This proves saturation to one.
\end{proof}
